% Generated draft: requires mathematical and source review.
\documentclass[11pt]{article}
\usepackage[margin=25mm]{geometry}
\usepackage{fontspec}
\setmainfont{Latin Modern Roman}
\usepackage{amsmath,amssymb,mathtools,unicode-math}
\setmathfont{Latin Modern Math}
\usepackage{enumitem,needspace,xurl,hyperref}
\setlength{\parindent}{0pt}
\setlength{\parskip}{5pt}
\setlength{\emergencystretch}{3em}
\newcommand{\N}{\mathbb{N}}
\newcommand{\Z}{\mathbb{Z}}
\newcommand{\Q}{\mathbb{Q}}
\newcommand{\R}{\mathbb{R}}
\newcommand{\C}{\mathbb{C}}
\newcommand{\F}{\mathbb{F}}
\newcommand{\A}{\mathbb{A}}
\newcommand{\PP}{\mathbb{P}}
\newcommand{\E}{\mathbb{E}}
\newcommand{\eps}{\varepsilon}
\newcommand{\dd}{\,\mathrm{d}}
\begin{document}
\section*{UnsolvedMath ID: 30000454}
\textbf{Source problem code:} catalogue entry
\subsection{Definitions and mathematical statement}
Let \(G\) be a finite group and let \(b\) be a block in the modular coefficient system intended in the source. In the usual setting, \(k\) has characteristic \(p\), \(\mathcal O\) is a suitable complete discrete valuation ring with residue field \(k\), and a block is a primitive central idempotent of the corresponding group algebra. The exact splitting and characteristic assumptions are not reproduced.

Let \(\operatorname{Out}(G)=\operatorname{Aut}(G)/\operatorname{Inn}(G)\), and let \(\operatorname{Out}(G)_b\) be the subgroup stabilizing the block. Let \(G_k(G,b)\) and \(LG_k(G,b)\) denote the source's global and local Grothendieck \(\mathcal O\)-modules, with their stated stabilizer actions. The exact global construction and, especially, the local construction are not supplied. A usual global modular Grothendieck-module convention is extension of scalars to \(\mathcal O\) of the Grothendieck group of finite-dimensional \(kGb\)-modules, but it must be matched to the source rather than substituted for an unspecified ordinary or twisted version.

Conjecture (C), as identified in the supplied verification, asks whether for every such \(G,b\) there is an isomorphism
\[
\Phi:G_k(G,b)\xrightarrow{\ \cong\ }LG_k(G,b)
\quad\text{of }\mathcal O[\operatorname{Out}(G)_b]\text{-modules}.
\]
Thus \(\Phi\) is an \(\mathcal O\)-linear bijection commuting with every stabilizer action. Equality of ranks, equality of ordinary characters after scalar extension, or an isomorphism over the fraction field alone is not the full integral module statement. No canonical choice of \(\Phi\) is required by the wording.

The original extraction used \(\operatorname{Aut}(G)_b\), while the cleaned, reportedly source-verified statement uses the outer stabilizer. The latter is retained. If inner automorphisms act trivially on both modules, the outer actions can be inflated to the automorphism stabilizer; that descent or triviality is a source property to verify, not an assumed correction without explanation.

The supplied assessment checked 2026-08-21 distinguishes a vanishing theorem in the surrounding passage from this conjectural global-to-local isomorphism and reports later development of the groups without a general proof for every block. Neither the vanishing theorem nor the precise local data are reproduced, so they are not invented or treated as resolving the conjecture.
\subsection{Short English statement}
Are the source's global and local block Grothendieck modules integrally isomorphic with their full outer-automorphism-stabilizer actions? The exact local construction and coefficient system are missing; agreement of ranks or scalar-extended characters would be weaker.
\subsection{Sources}
\begin{itemize}
\item Local UnsolvedMath registry, permanent ID 30000454, catalogue entry, \emph{Integral Local Grothendieck-Module Isomorphism}.
\item \emph{Representations of Finite Groups} (2007), Oberwolfach Reports, pp. 895--978, DOI \url{https://doi.org/10.4171/owr/2006/15}, conjecture (C) as identified by supplied verification.
\item Puig, supplied label ``ordinary Grothendieck groups'': \url{https://arxiv.org/abs/1004.1615}.
\item Puig, supplied label ``categorizations of Grothendieck groups'': \url{https://arxiv.org/abs/1511.05033}.
\item Supplied literature triage checked 2026-08-21; the vanishing theorem and detailed group constructions are not reproduced.
\end{itemize}
\end{document}
